\section{Theoretical background}
\subsection{Fracture mechanics}
[...]

The application of fracture mechanics to structures subjected to fatigue was first introduced in the early 1960s by \textcite{paris1997rational}, who demonstrated that the stress intensity factor, $K$, could be used to characterize crack growth driven by cyclic loading. Because $K$ is directly proportional to both the applied load and the component's geometry, variations in either parameter dictate the stress field at the crack tip. Assuming the crack length and overall geometry remain essentially constant during a single load cycle, the load ratio, $R$, can be expressed equivalently as the ratio of the minimum and maximum stress intensity factors, independent of geometry:

\begin{equation}
	R = \frac{P_{max}}{P_{min}} = \frac{K_{max}}{K_{min}}.
\end{equation}

Although the elastic singularity mathematically governs the stress field near the crack tip, localized plasticity is the dominant mechanism in the immediate tip region. However, if this plastic zone is sufficiently small, $K$ alone can represent the crack tip stress state. This fundamental assumption is referred to as the principle of similitude. Since $R$ and $K$ can represent stress state variation and intensity, any phenomena occurring close to the crack tip is governed by those two parameters such as the crack length variation per cycle:

\begin{equation}
	\frac{\Delta a}{\Delta N} = f(K_{max}, R).
\end{equation} 
